% !TEX root = chapter-standalone.tex


\section{Monotone relations}
\label{sec:monotone-relations}

In \cref{chap:relation} we generalized functions to relations: every function $\mapa \colon \setA \sto \setB$ gives a relation, namely its graph, but not every relation is the graph of a function.
In \cref{chap:monotone-functions} we then brought order into the picture, and studied \SY{monotone functions} between \SY{preorders}.
In this chapter we do the same for relations: we study relations between \SY{preorders} that are compatible with the preorder relations on their source and target.
We call them \emph{monotone relations}.

Two questions will guide us.
First, what is the right notion of compatibility between a relation and the preorders at hand?
Second, in what sense do monotone relations generalize monotone functions, in the same way that relations generalize functions?
The answer to the second question is more subtle than one might expect.
As we will see, the graph of a monotone function is in general \emph{not} a monotone relation, and we will need a small twist.

\begin{remark}
    In the next chapter (\cref{chap:application-design-problems}), monotone relations will appear prominently in co-design, where they describe which combinations of functionality and requirements are feasible.
    In this chapter we develop the general theory, independently of that application.
\end{remark}

\subsection{Definition and first examples}

\begin{ctdefinition}[Monotone relation]
    \label{def:monotone-relation}
    Let $\posAdefinition$ and $\posBdefinition$ be \SY{preorders}.
    
    A \maindef{monotone relation} $\monrelM \colon \posA \profto \posB$ is:

    \constit
    \begin{enumerate}
        \item a \SY{relation} $\monrelM \setsubseteq \posAset \cartprod \posBset$ (\cref{def:binary-relation});
    \end{enumerate}
    \condit
    \begin{enumerate}
        \item $\monrelM$ is \emph{closed downward in the source}: for all $\posAel, \posAel' \setin \posAset$ and $\posBel \setin \posBset$,
              \begin{equation}\label{eq:monrel-cond-source}
                  \prfperiod{
                      \posAel' \posAleq \posAel
                  }{
                      \inrel{\posAel}{\monrelM}{\posBel}
                  }{
                      \inrel{\posAel'}{\monrelM}{\posBel}
                  }
              \end{equation}
        \item $\monrelM$ is \emph{closed upward in the target}: for all $\posAel \setin \posAset$ and $\posBel, \posBel' \setin \posBset$,
              \begin{equation}\label{eq:monrel-cond-target}
                  \prfperiod{
                      \inrel{\posAel}{\monrelM}{\posBel}
                  }{
                      \posBel \posBleq \posBel'
                  }{
                      \inrel{\posAel}{\monrelM}{\posBel'}
                  }
              \end{equation}
    \end{enumerate}
\end{ctdefinition}

In words: if $\posAel$ is related to $\posBel$, then everything below $\posAel$ is related to $\posBel$, and $\posAel$ is related to everything above $\posBel$.
We use the slashed arrow ``$\profto$'' for monotone relations, to distinguish them from \SY{monotone functions}, for which we use ``$\mto$''.
We will usually name monotone relations $\monrelM$ and $\monrelN$ (with primes or numbered subscripts, such as $\monrelM'$ or $\monrelM_1$, when we need more), and keep the letters $\relA$ and $\relB$ for relations that are not assumed to be monotone.

The two conditions of \cref{def:monotone-relation} can be packaged into a single one.

\begin{lemma}
    \label{lem:monrel-single-condition}
    Let \posA and \posB be \SY{preorders}.
    A relation $\relA \setsubseteq \posAset \cartprod \posBset$ is a monotone relation $\posA \profto \posB$ if and only if, for all $\posAel, \posAel' \setin \posAset$ and $\posBel, \posBel' \setin \posBset$,
    \begin{equation}\label{eq:monrel-single-condition}
        \prfperiod{
            \posAel' \posAleq \posAel
        }{
            \inrel{\posAel}{\relA}{\posBel}
        }{
            \posBel \posBleq \posBel'
        }{
            \inrel{\posAel'}{\relA}{\posBel'}
        }
    \end{equation}
\end{lemma}
\begin{proof}
    Suppose first that \relA is a monotone relation, and that $\posAel' \posAleq \posAel$, $\inrel{\posAel}{\relA}{\posBel}$, and $\posBel \posBleq \posBel'$.
    By \cref{eq:monrel-cond-source}, from $\posAel' \posAleq \posAel$ and $\inrel{\posAel}{\relA}{\posBel}$ we get $\inrel{\posAel'}{\relA}{\posBel}$.
    Then, by \cref{eq:monrel-cond-target}, from $\inrel{\posAel'}{\relA}{\posBel}$ and $\posBel \posBleq \posBel'$ we get $\inrel{\posAel'}{\relA}{\posBel'}$.

    Conversely, suppose that \relA satisfies \cref{eq:monrel-single-condition}.
    To obtain \cref{eq:monrel-cond-source}, suppose $\posAel' \posAleq \posAel$ and $\inrel{\posAel}{\relA}{\posBel}$.
    By reflexivity, $\posBel \posBleq \posBel$, so \cref{eq:monrel-single-condition} with $\posBel' = \posBel$ gives $\inrel{\posAel'}{\relA}{\posBel}$.
    To obtain \cref{eq:monrel-cond-target}, suppose $\inrel{\posAel}{\relA}{\posBel}$ and $\posBel \posBleq \posBel'$.
    By reflexivity, $\posAel \posAleq \posAel$, so \cref{eq:monrel-single-condition} with $\posAel' = \posAel$ gives $\inrel{\posAel}{\relA}{\posBel'}$.
\end{proof}

\begin{example}[The preorder relation itself]
    \label{exa:monrel-leq}
    Let $\posAdefinition$ be a \SY{preorder}.
    The relation $\posAleq \setsubseteq \posAset \cartprod \posAset$, seen as a relation from \posA to \posA, is a monotone relation
    \begin{equation}
        \posAleq \colon \posA \profto \posA.
    \end{equation}
    Indeed, if $\posAel' \posAleq \posAel$ and $\posAel \posAleq \posAel''$, then $\posAel' \posAleq \posAel''$ by transitivity; this gives \cref{eq:monrel-cond-source}.
    Similarly, if $\posAel \posAleq \posAel''$ and $\posAel'' \posAleq \posAel'''$, then $\posAel \posAleq \posAel'''$ by transitivity; this gives \cref{eq:monrel-cond-target}.
    We will see below (\cref{lem:monrel-leq-neutral}) that this monotone relation plays the role of an identity.
\end{example}

\begin{example}[Relations between sets]
    \label{exa:monrel-discrete}
    Let \setA and \setB be sets, and view them as discrete posets $\tup{\setA, {{=}}}$ and $\tup{\setB, {{=}}}$ (\cref{ex:discreteposet}).
    Then \emph{every} relation $\relA \setsubseteq \setA \cartprod \setB$ is a monotone relation $\tup{\setA, {{=}}} \profto \tup{\setB, {{=}}}$.
    Indeed, in \cref{eq:monrel-cond-source} the hypothesis $\posAel' = \posAel$ makes the conclusion identical to the hypothesis $\inrel{\posAel}{\relA}{\posBel}$, and similarly for \cref{eq:monrel-cond-target}.

    This parallels \cref{lem:discrete-is-monotone}, which says that every function out of a discrete poset is monotone.
    (For monotone relations, we need \emph{both} source and target to be discrete, because both conditions of \cref{def:monotone-relation} involve an order.)
    So ordinary relations between sets are the special case of monotone relations in which there is no order to respect.
\end{example}

\begin{example}[Empty and full relations]
    For any \SY{preorders} \posA and \posB, the empty relation $\Emptyset \setsubseteq \posAset \cartprod \posBset$ and the full relation $\posAset \cartprod \posBset$ are monotone relations $\posA \profto \posB$: for the empty relation the hypotheses of both conditions are never satisfied, and for the full relation the conclusions are always satisfied.
\end{example}

\begin{example}[Paying with coins]
    \label{exa:monrel-coins}
    Let $\posA = \natswithleq$, whose elements we think of as amounts of money, in francs.
    Let $\posB = \natswithleq \Ptimes \natswithleq$ be the product preorder (\cref{def:poset-product}), whose elements $\tup{\elnc{1}, \elnc{2}}$ we think of as a number $\elnc{1}$ of one-franc coins together with a number $\elnc{2}$ of two-franc coins.
    Define the relation $\monrelM \colon \posA \profto \posB$ by
    \begin{equation}
        \inrel{\ela}{\monrelM}{\tup{\elnc{1}, \elnc{2}}}
        \quad \text{if and only if} \quad
        \ela \leq \elnc{1} + 2 \elnc{2},
    \end{equation}
    that is, if the coins suffice to pay the amount $\ela$ (possibly overpaying).
    This is a monotone relation.
    If the coins suffice to pay $\ela$, they also suffice to pay any smaller amount $\ela' \leq \ela$, since $\ela' \leq \ela \leq \elnc{1} + 2 \elnc{2}$.
    And if we have at least as many coins of each kind, that is, $\tup{\elnc{1}, \elnc{2}} \posBleq \tup{\elnc{1}', \elnc{2}'}$, then $\elnc{1} + 2\elnc{2} \leq \elnc{1}' + 2\elnc{2}'$, so the new coins suffice as well.
    We will come back to this example several times.
\end{example}

\begin{example}[Course admission]
    \label{exa:monrel-courses}
    This example involves \SY{preorders} that are not \SY{posets}.
    Let \setA be a set of courses, each with a minimum admission grade given by a function $\stylemaps{req} \colon \setA \sto \reals$.
    Let \setB be a set of students, each with an average grade given by a function $\stylemaps{avg} \colon \setB \sto \reals$.
    We preorder the courses by their requirements and the students by their grades:
    \begin{equation}
        \elna{1} \posleqof{\posA} \elna{2} \ \Eqv\ \stylemaps{req}(\elna{1}) \leq \stylemaps{req}(\elna{2}),
        \qquad
        \elnb{1} \posleqof{\posB} \elnb{2} \ \Eqv\ \stylemaps{avg}(\elnb{1}) \leq \stylemaps{avg}(\elnb{2}).
    \end{equation}
    These relations are reflexive and transitive, because $\leq$ on $\reals$ is.
    But they are not antisymmetric as soon as two different courses have the same requirement, or two different students have the same average.
    Write $\posA = \tup{\setA, \posleqof{\posA}}$ and $\posB = \tup{\setB, \posleqof{\posB}}$ for the resulting \SY{preorders}.
    The relation ``course $\ela$ is open to student $\elb$'',
    \begin{equation}
        \inrel{\ela}{\monrelM}{\elb}
        \quad \text{if and only if} \quad
        \stylemaps{req}(\ela) \leq \stylemaps{avg}(\elb),
    \end{equation}
    is a monotone relation $\monrelM \colon \posA \profto \posB$.
    If a course is open to a student, then so is every course with a lower or equal requirement, and the course is also open to every student with a higher or equal average.
\end{example}

\begin{example}[A graph that is not a monotone relation]
    \label{exa:graph-not-monrel}
    The identity function on $\natswithleq$ is monotone, and its graph is the relation
    \begin{equation}
        \makeset{\tup{\ela, \ela} \mid \ela \setin \natnumbers},
    \end{equation}
    that is, equality.
    This is \emph{not} a monotone relation $\natswithleq \profto \natswithleq$: we have $\tup{1, 1}$ in the graph and $1 \leq 2$, but $\tup{1, 2}$ is not in the graph, so \cref{eq:monrel-cond-target} fails.
    More generally, the graph of a monotone function is almost never closed upward in the target.
    So we cannot view monotone functions as monotone relations simply by taking graphs, as we did for functions in \cref{def:functions_as_relations}.
    We will see how to fix this in \cref{sec:monrel-from-monotone-functions}.
\end{example}

\subsection{Monotone relations as truth-valued matrices}

In \cref{rem:rel-three-fun-descriptions} we saw that a relation $\relA \colon \setA \sto \setB$ between sets can be described equivalently by its indicator function $\phi_\relA \colon \setA \cartprod \setB \sto \boolset$.
When \setA and \setB are finite, we can write $\phi_\relA$ as a matrix with entries in $\boolset$, with rows indexed by \setA and columns indexed by \setB.

The same can be done for monotone relations; however, in the world of \SY{preorders}, we need to introduce a small twist.
Monotonicity of \relA does not correspond to monotonicity of $\phi_\relA$ on $\posA \Ptimes \posB$, but on $\posAop \Ptimes \posB$: the order on the source has to be reversed.

\begin{lemma}
    \label{lem:monrel-as-bool-map}
    Let \posA and \posB be \SY{preorders}, let $\relA \setsubseteq \posAset \cartprod \posBset$ be a relation, and let $\phi_\relA \colon \posAset \cartprod \posBset \sto \boolset$ be its indicator function.
    The following are equivalent:
    \begin{enumerate}
        \item $\relA$ is a monotone relation $\posA \profto \posB$;
        \item $\relA$ is an \SY{upper set} (\cref{def:upperset}) of the product preorder $\posAop \Ptimes \posB$ (\cref{def:poset-opposite,def:poset-product});
        \item $\phi_\relA$ is a \SY{monotone function} $\posAop \Ptimes \posB \mto \Bool$ (\cref{def:bool-poset}).
    \end{enumerate}
\end{lemma}
% Proof commented out: it overlaps with the graded exercise DPsAsUpperSets in the design-problems chapter.
% \begin{proof}
%     We first unwind the order of $\posAop \Ptimes \posB$.
%     By the definitions of the product and of the opposite, for all $\tup{\posAel, \posBel}, \tup{\posAel', \posBel'} \setin \posAset \cartprod \posBset$ we have
%     \begin{equation}\label{eq:order-Aop-times-B}
%         \tup{\posAel, \posBel} \posleqof{\posAop \Ptimes \posB} \tup{\posAel', \posBel'}
%         \quad \text{if and only if} \quad
%         \posAel' \posAleq \posAel
%         \ \text{ and } \
%         \posBel \posBleq \posBel'.
%     \end{equation}
% 
%     \emph{(1)$\Eqv$(2).}
%     By \cref{def:upperset} and \cref{eq:order-Aop-times-B}, $\relA$ is an upper set of $\posAop \Ptimes \posB$ if and only if, whenever $\inrel{\posAel}{\relA}{\posBel}$, $\posAel' \posAleq \posAel$, and $\posBel \posBleq \posBel'$, we have $\inrel{\posAel'}{\relA}{\posBel'}$.
%     This is exactly condition \cref{eq:monrel-single-condition}, which by \cref{lem:monrel-single-condition} is equivalent to \relA being a monotone relation.
% 
%     \emph{(2)$\Eqv$(3).}
%     Write $\posC = \posAop \Ptimes \posB$.
%     A function $\phi \colon \posCset \sto \boolset$ fails to be monotone $\posC \mto \Bool$ exactly when there are $\elc \posCleq \elc'$ such that $\phi(\elc) \posleqof\Bool \phi(\elc')$ does not hold.
%     In $\Bool$, the only pair of elements for which $\posleqof{\Bool}$ does not hold is $\tup{\true, \false}$.
%     So $\phi$ is monotone if and only if, whenever $\elc \posCleq \elc'$ and $\phi(\elc) = \true$, also $\phi(\elc') = \true$.
%     For $\phi = \phi_\relA$, this says precisely that $\relA = \makeset{\elc \setin \posCset \mid \phi_\relA(\elc) = \true}$ is an upper set of \posC.
% \end{proof}

\begin{remark}[Why the twist?]
    It is instructive to see what goes wrong without the twist.
    A relation whose indicator function is monotone on $\posA \Ptimes \posB$ is an upper set of $\posA \Ptimes \posB$: it is closed upward in \emph{both} the source and the target.
    The preorder relation $\posAleq$ of \cref{exa:monrel-leq} does not have this property in general: from $\posAel \posAleq \posAel''$ and $\posAel \posAleq \posAel'$ we cannot conclude that $\posAel' \posAleq \posAel''$.
    Reversing the order on the source is exactly what makes the preorder relation itself, and more generally the constructions of \cref{sec:monrel-from-monotone-functions}, into monotone relations.
\end{remark}

\begin{example}[A monotone relation as a matrix]
    \label{exa:monrel-matrix}
    Let $\posA = \posB = \tup{\makeset{1, 2, 3}, \leq}$, and consider the monotone relation $\leq \colon \posA \profto \posB$ of \cref{exa:monrel-leq}.
    We write its indicator function as a matrix, with $1$ for $\true$ and $0$ for $\false$.
    We list the rows in increasing order for $\posAop$ (that is, in \emph{decreasing} order for \posA), and the columns in increasing order for \posB:
    \begin{equation}
        \begin{array}{c|ccc}
              & 1 & 2 & 3 \\ \hline
            3 & 0 & 0 & 1 \\
            2 & 0 & 1 & 1 \\
            1 & 1 & 1 & 1
        \end{array}
    \end{equation}
    Monotonicity is visible in the shape of the matrix: whenever an entry is $1$, so are all entries to its right (closure upward in the target) and all entries below it (closure downward in the source).
    For relations between finite totally ordered sets, listed in this way, this ``staircase'' shape characterizes the monotone relations.
\end{example}

\begin{example}[Upper sets and lower sets]
    \label{exa:monrel-upper-lower-sets}
    Let $\singleton = \makeset{\singletonel}$ be the one-element preorder, with $\singletonel \posleq \singletonel$.
    A monotone relation $\monrelM \colon \singleton \profto \posB$ is the same thing as an \SY{upper set} of \posB.
    Indeed, such a relation is determined by the subset $\stylesets{U} = \makeset{\posBel \setin \posBset \mid \inrel{\singletonel}{\monrelM}{\posBel}}$.
    Condition \cref{eq:monrel-cond-source} holds automatically, because the only element of the source is $\singletonel$, and condition \cref{eq:monrel-cond-target} says precisely that $\stylesets{U}$ is an upper set.

    Dually, a monotone relation $\monrelM \colon \posA \profto \singleton$ is the same thing as a \SY{lower set} of \posA, namely $\makeset{\posAel \setin \posAset \mid \inrel{\posAel}{\monrelM}{\singletonel}}$.
    So the upper and lower sets of \cref{sec:UpperLowerSets} are special cases of monotone relations.
\end{example}

\subsection{Composing monotone relations}

Monotone relations can be composed using the composition of relations (\cref{def:composition-of-relations}), and the result is again a monotone relation.

\begin{lemma}
    \label{lem:monrel-composition}
    Let \posA, \posB, and \posC be \SY{preorders}, and let $\monrelM \colon \posA \profto \posB$ and $\monrelN \colon \posB \profto \posC$ be monotone relations.
    Then the composite relation $\monrelM \mthen \monrelN$ is a monotone relation $\posA \profto \posC$.
\end{lemma}
\begin{proof}
    Recall that $\inrel{\posAel}{(\monrelM \mthen \monrelN)}{\posCel}$ means that there exists $\posBel \setin \posBset$ with $\inrel{\posAel}{\monrelM}{\posBel}$ and $\inrel{\posBel}{\monrelN}{\posCel}$.
    \begin{enumerate}
        \item Closure downward in the source: suppose $\posAel' \posAleq \posAel$ and $\inrel{\posAel}{(\monrelM \mthen \monrelN)}{\posCel}$.
              Choose $\posBel$ with $\inrel{\posAel}{\monrelM}{\posBel}$ and $\inrel{\posBel}{\monrelN}{\posCel}$.
              Since \monrelM is closed downward in the source, $\inrel{\posAel'}{\monrelM}{\posBel}$.
              Together with $\inrel{\posBel}{\monrelN}{\posCel}$, this shows $\inrel{\posAel'}{(\monrelM \mthen \monrelN)}{\posCel}$.
        \item Closure upward in the target: suppose $\inrel{\posAel}{(\monrelM \mthen \monrelN)}{\posCel}$ and $\posCel \posCleq \posCel'$.
              Choose $\posBel$ with $\inrel{\posAel}{\monrelM}{\posBel}$ and $\inrel{\posBel}{\monrelN}{\posCel}$.
              Since \monrelN is closed upward in the target, $\inrel{\posBel}{\monrelN}{\posCel'}$.
              Together with $\inrel{\posAel}{\monrelM}{\posBel}$, this shows $\inrel{\posAel}{(\monrelM \mthen \monrelN)}{\posCel'}$.
    \end{enumerate}
\end{proof}

What about identities?
For sets, the identity relation on \setA is equality (\cref{def:identity-relation}).
But equality on \posAset is in general \emph{not} a monotone relation $\posA \profto \posA$, as \cref{exa:graph-not-monrel} shows: it is a monotone relation only if $\posAleq$ is itself equality.
Instead, the role of the identity is played by the preorder relation $\posAleq$ of \cref{exa:monrel-leq}.

\begin{lemma}
    \label{lem:monrel-leq-neutral}
    Let \posA and \posB be \SY{preorders}, and let $\monrelM \colon \posA \profto \posB$ be a monotone relation.
    Then
    \begin{equation}
        \posAleq \mthen \monrelM = \monrelM
        \qqand
        \monrelM \mthen \posBleq = \monrelM.
    \end{equation}
\end{lemma}
\begin{proof}
    We prove the first equation; the second one is proved in the same way, using closure upward in the target instead of closure downward in the source.
    \begin{itemize}
        \item $\posAleq \mthen \monrelM \setsubseteq \monrelM$: suppose $\inrel{\posAel'}{(\posAleq \mthen \monrelM)}{\posBel}$.
              Then there is $\posAel \setin \posAset$ with $\posAel' \posAleq \posAel$ and $\inrel{\posAel}{\monrelM}{\posBel}$.
              Since \monrelM is closed downward in the source, $\inrel{\posAel'}{\monrelM}{\posBel}$.
        \item $\monrelM \setsubseteq \posAleq \mthen \monrelM$: suppose $\inrel{\posAel}{\monrelM}{\posBel}$.
              By reflexivity, $\posAel \posAleq \posAel$, so $\posAel$ itself witnesses that $\inrel{\posAel}{(\posAleq \mthen \monrelM)}{\posBel}$.
    \end{itemize}
\end{proof}

\begin{remark}
    Taking $\monrelM = \posAleq$ in \cref{lem:monrel-leq-neutral}, we get $\posAleq \mthen \posAleq = \posAleq$.
    Here the inclusion $\setsubseteq$ is transitivity, and the inclusion $\setsupseteq$ follows from reflexivity.
    So the two axioms of a \SY{preorder} are exactly what makes its preorder relation behave like an identity for composition.
\end{remark}

We can also rephrase \cref{def:monotone-relation} entirely in the language of composition of relations.

\begin{lemma}
    \label{lem:monrel-relational}
    Let \posA and \posB be \SY{preorders}, and let $\relA \setsubseteq \posAset \cartprod \posBset$ be a relation.
    Then \relA is a monotone relation $\posA \profto \posB$ if and only if
    \begin{equation}\label{eq:monrel-relational}
        (\posAleq \mthen \relA) \mthen \posBleq \ \setsubseteq\ \relA,
    \end{equation}
    and in that case the two sides of \cref{eq:monrel-relational} are equal.
\end{lemma}
\begin{proof}
    Unwinding the definition of composition twice, $\inrel{\posAel'}{((\posAleq \mthen \relA) \mthen \posBleq)}{\posBel'}$ holds if and only if there exist $\posAel \setin \posAset$ and $\posBel \setin \posBset$ with $\posAel' \posAleq \posAel$, $\inrel{\posAel}{\relA}{\posBel}$, and $\posBel \posBleq \posBel'$.
    So \cref{eq:monrel-relational} says exactly that condition \cref{eq:monrel-single-condition} holds, and by \cref{lem:monrel-single-condition} this is equivalent to \relA being monotone.
    Moreover, the reverse inclusion always holds: if $\inrel{\posAel}{\relA}{\posBel}$, then $\posAel \posAleq \posAel$ and $\posBel \posBleq \posBel$ by reflexivity, so $\inrel{\posAel}{((\posAleq \mthen \relA) \mthen \posBleq)}{\posBel}$.
\end{proof}

\subsection{From monotone functions to monotone relations}
\label{sec:monrel-from-monotone-functions}

We now come back to the question of how monotone functions can be seen as monotone relations.
\Cref{exa:graph-not-monrel} shows that the graph of a monotone function $\mapa \colon \posA \mto \posB$ is in general not closed upward in the target.
The natural fix is to close it upward: instead of relating $\posAel$ only to $\mapa(\posAel)$, we relate $\posAel$ to everything above $\mapa(\posAel)$.
There is also a second, dual, way of turning $\mapa$ into a monotone relation, which goes in the opposite direction.

\begin{definition}[Companion and conjoint]
    \label{def:monrel-companion-conjoint}
    Let \posA and \posB be \SY{preorders}, and let $\mapa \colon \posA \mto \posB$ be a \SY{monotone function}.
    \begin{enumerate}
        \item The \emph{companion} of $\mapa$ is the relation $\comp{\mapa} \colon \posA \profto \posB$ given by
              \begin{equation}
                  \inrel{\posAel}{\comp{\mapa}}{\posBel}
                  \quad \text{if and only if} \quad
                  \mapa(\posAel) \posBleq \posBel.
              \end{equation}
        \item The \emph{conjoint} of $\mapa$ is the relation $\conj{\mapa} \colon \posB \profto \posA$ given by
              \begin{equation}
                  \inrel{\posBel}{\conj{\mapa}}{\posAel}
                  \quad \text{if and only if} \quad
                  \posBel \posBleq \mapa(\posAel).
              \end{equation}
    \end{enumerate}
\end{definition}

Note that the conjoint goes in the \emph{opposite} direction to $\mapa$, from \posB to \posA.
We will meet these two constructions again, in the setting of design problems, in \cref{def:comp_conj}.

\begin{lemma}
    \label{lem:companion-conjoint-monotone}
    Let $\mapa \colon \posA \mto \posB$ be a \SY{monotone function} between \SY{preorders}.
    Then $\comp{\mapa} \colon \posA \profto \posB$ and $\conj{\mapa} \colon \posB \profto \posA$ are monotone relations.
\end{lemma}
% Proof commented out: it overlaps with the graded exercise DPsFromMonotoneMaps in the design-problems chapter.
% \begin{proof}
%     For the companion:
%     \begin{enumerate}
%         \item Closure downward in the source: suppose $\posAel' \posAleq \posAel$ and $\mapa(\posAel) \posBleq \posBel$.
%               Since $\mapa$ is monotone, $\mapa(\posAel') \posBleq \mapa(\posAel)$, and by transitivity $\mapa(\posAel') \posBleq \posBel$.
%         \item Closure upward in the target: suppose $\mapa(\posAel) \posBleq \posBel$ and $\posBel \posBleq \posBel'$.
%               By transitivity, $\mapa(\posAel) \posBleq \posBel'$.
%     \end{enumerate}
%     For the conjoint (whose source is \posB and whose target is \posA):
%     \begin{enumerate}
%         \item Closure downward in the source: suppose $\posBel' \posBleq \posBel$ and $\posBel \posBleq \mapa(\posAel)$.
%               By transitivity, $\posBel' \posBleq \mapa(\posAel)$.
%         \item Closure upward in the target: suppose $\posBel \posBleq \mapa(\posAel)$ and $\posAel \posAleq \posAel'$.
%               Since $\mapa$ is monotone, $\mapa(\posAel) \posBleq \mapa(\posAel')$, and by transitivity $\posBel \posBleq \mapa(\posAel')$.
%     \end{enumerate}
% \end{proof}

% Commented out together with the proof above (it describes that proof):
% Notice that the monotonicity of $\mapa$ is used exactly once in each case: for the companion in the source condition, and for the conjoint in the target condition.

\begin{example}
    \label{exa:companion-identity}
    For the identity function $\mapidat{\posAset} \colon \posA \mto \posA$, both the companion and the conjoint are the preorder relation:
    \begin{equation}
        \comp{\mapidat{\posAset}} = \posAleq = \conj{\mapidat{\posAset}}.
    \end{equation}
    Indeed, $\inrel{\posAel}{\comp{\mapidat{\posAset}}}{\posAel'}$ means $\posAel \posAleq \posAel'$, and $\inrel{\posAel}{\conj{\mapidat{\posAset}}}{\posAel'}$ also means $\posAel \posAleq \posAel'$.
    So the identity monotone function corresponds to the identity monotone relation of \cref{lem:monrel-leq-neutral}.
\end{example}

The companion and the conjoint are related to the graph of $\mapa$, and to the transpose of relations (\cref{def:relation-transpose}), as follows.
Write $\aword{graph}(\mapa) = \makeset{\tup{\posAel, \mapa(\posAel)} \mid \posAel \setin \posAset}$ for the graph of $\mapa$.

\begin{lemma}
    \label{lem:companion-via-graph}
    Let $\mapa \colon \posA \mto \posB$ be a \SY{monotone function} between \SY{preorders}.
    Then
    \begin{equation}
        \comp{\mapa} = \aword{graph}(\mapa) \mthen \posBleq
        \qqand
        \conj{\mapa} = \posBleq \mthen \aword{graph}(\mapa)\reltransp.
    \end{equation}
    Moreover, $\comp{\mapa}$ is the smallest monotone relation $\posA \profto \posB$ that includes $\aword{graph}(\mapa)$.
\end{lemma}
\begin{proof}
    For the first equation: $\inrel{\posAel}{(\aword{graph}(\mapa) \mthen \posBleq)}{\posBel}$ if and only if there is $\posBel_0 \setin \posBset$ with $\posBel_0 = \mapa(\posAel)$ and $\posBel_0 \posBleq \posBel$, that is, if and only if $\mapa(\posAel) \posBleq \posBel$.
    For the second equation: $\inrel{\posBel}{(\posBleq \mthen \aword{graph}(\mapa)\reltransp)}{\posAel}$ if and only if there is $\posBel_0 \setin \posBset$ with $\posBel \posBleq \posBel_0$ and $\posBel_0 = \mapa(\posAel)$, that is, if and only if $\posBel \posBleq \mapa(\posAel)$.

    For the last claim, note first that $\aword{graph}(\mapa) \setsubseteq \comp{\mapa}$, since $\mapa(\posAel) \posBleq \mapa(\posAel)$ by reflexivity.
    Now let $\monrelM \colon \posA \profto \posB$ be any monotone relation with $\aword{graph}(\mapa) \setsubseteq \monrelM$, and suppose $\inrel{\posAel}{\comp{\mapa}}{\posBel}$, that is, $\mapa(\posAel) \posBleq \posBel$.
    Since $\tup{\posAel, \mapa(\posAel)} \setin \aword{graph}(\mapa) \setsubseteq \monrelM$ and \monrelM is closed upward in the target, we get $\inrel{\posAel}{\monrelM}{\posBel}$.
    Hence $\comp{\mapa} \setsubseteq \monrelM$.
\end{proof}

This lemma is the precise sense in which the companion ``repairs'' the graph: it is the graph, closed upward in the target.

\subsubsection{Monotone relations generalize monotone functions}

For sets, a function can be recovered from its graph.
For monotone functions, we can recover a function from its companion, up to the ambiguity that a \SY{preorder} cannot see.

\begin{lemma}
    \label{lem:companion-determines-function}
    Let $\mapa, \mapb \colon \posA \mto \posB$ be \SY{monotone functions} between \SY{preorders}.
    Then $\comp{\mapa} = \comp{\mapb}$ if and only if, for all $\posAel \setin \posAset$,
    \begin{equation}
        \mapa(\posAel) \posBleq \mapb(\posAel)
        \qqand
        \mapb(\posAel) \posBleq \mapa(\posAel).
    \end{equation}
    In particular, if \posB is a \SY{poset}, then $\comp{\mapa} = \comp{\mapb}$ if and only if $\mapa = \mapb$.
    The same statements hold for conjoints.
\end{lemma}
\begin{proof}
    Suppose $\comp{\mapa} = \comp{\mapb}$, and let $\posAel \setin \posAset$.
    By reflexivity, $\mapb(\posAel) \posBleq \mapb(\posAel)$, that is, $\inrel{\posAel}{\comp{\mapb}}{\mapb(\posAel)}$.
    Hence $\inrel{\posAel}{\comp{\mapa}}{\mapb(\posAel)}$, that is, $\mapa(\posAel) \posBleq \mapb(\posAel)$.
    Exchanging the roles of $\mapa$ and $\mapb$ gives $\mapb(\posAel) \posBleq \mapa(\posAel)$.

    Conversely, suppose that the two inequalities hold for all $\posAel$, and suppose $\inrel{\posAel}{\comp{\mapa}}{\posBel}$, that is, $\mapa(\posAel) \posBleq \posBel$.
    Then $\mapb(\posAel) \posBleq \mapa(\posAel) \posBleq \posBel$, so $\inrel{\posAel}{\comp{\mapb}}{\posBel}$ by transitivity.
    This shows $\comp{\mapa} \setsubseteq \comp{\mapb}$, and the reverse inclusion is obtained by exchanging the roles of $\mapa$ and $\mapb$.

    If \posB is a \SY{poset}, the two inequalities together are equivalent to $\mapa(\posAel) = \mapb(\posAel)$ by antisymmetry, so they hold for all $\posAel$ if and only if $\mapa = \mapb$.
    For conjoints, the argument is the same, using $\inrel{\mapb(\posAel)}{\conj{\mapb}}{\posAel}$ in place of $\inrel{\posAel}{\comp{\mapb}}{\mapb(\posAel)}$.
\end{proof}

So, when the target is a \SY{poset}, taking companions gives an injective map from monotone functions $\posA \mto \posB$ to monotone relations $\posA \profto \posB$.
In this sense we can view monotone functions as special monotone relations, just as we viewed functions as special relations via their graphs (\cref{def:functions_as_relations}).
When the target is only a \SY{preorder}, two monotone functions have the same companion exactly when their values are equivalent in \posB at every input; the companion forgets precisely the information that the preorder cannot distinguish.

Which monotone relations are companions?
For a monotone relation $\monrelM \colon \posA \profto \posB$ and $\posAel \setin \posAset$, write
\begin{equation}
    \hat{\monrelM}(\posAel) = \makeset{\posBel \setin \posBset \mid \inrel{\posAel}{\monrelM}{\posBel}}
\end{equation}
for the set of elements to which $\posAel$ is related, as in \cref{rem:rel-three-fun-descriptions}.

\begin{lemma}
    \label{lem:which-monrels-are-companions}
    Let $\monrelM \colon \posA \profto \posB$ be a monotone relation between \SY{preorders}.
    Then $\monrelM = \comp{\mapa}$ for some \SY{monotone function} $\mapa \colon \posA \mto \posB$ if and only if, for every $\posAel \setin \posAset$, the set $\hat{\monrelM}(\posAel)$ has a least element, that is, an element $\posBel_0 \setin \hat{\monrelM}(\posAel)$ with $\posBel_0 \posBleq \posBel$ for all $\posBel \setin \hat{\monrelM}(\posAel)$.
\end{lemma}
\begin{proof}
    Suppose $\monrelM = \comp{\mapa}$.
    Then $\hat{\monrelM}(\posAel) = \makeset{\posBel \mid \mapa(\posAel) \posBleq \posBel}$, which contains $\mapa(\posAel)$ by reflexivity; and $\mapa(\posAel)$ is below every element of $\hat{\monrelM}(\posAel)$ by definition.
    So $\mapa(\posAel)$ is a least element.

    Conversely, suppose that each $\hat{\monrelM}(\posAel)$ has a least element, and choose one for each $\posAel$; call it $\mapa(\posAel)$.
    (In a \SY{poset} the least element is unique, so there is no choice to make.)
    This defines a function $\mapa \colon \posAset \sto \posBset$.
    \begin{itemize}
        \item $\mapa$ is monotone: suppose $\posAel \posAleq \posAel'$.
              We have $\inrel{\posAel'}{\monrelM}{\mapa(\posAel')}$, so, since \monrelM is closed downward in the source, also $\inrel{\posAel}{\monrelM}{\mapa(\posAel')}$, that is, $\mapa(\posAel') \setin \hat{\monrelM}(\posAel)$.
              Since $\mapa(\posAel)$ is least in $\hat{\monrelM}(\posAel)$, we get $\mapa(\posAel) \posBleq \mapa(\posAel')$.
        \item $\comp{\mapa} = \monrelM$: if $\inrel{\posAel}{\monrelM}{\posBel}$, then $\posBel \setin \hat{\monrelM}(\posAel)$, so $\mapa(\posAel) \posBleq \posBel$ because $\mapa(\posAel)$ is least.
              Conversely, if $\mapa(\posAel) \posBleq \posBel$, then from $\inrel{\posAel}{\monrelM}{\mapa(\posAel)}$ and closure upward in the target we get $\inrel{\posAel}{\monrelM}{\posBel}$.
    \end{itemize}
\end{proof}

\begin{example}[Monotone relations that are not companions]
    The monotone relation of \cref{exa:monrel-coins} is not the companion of any monotone function.
    The amount $3$ can be paid with three one-franc coins, or with one coin of each kind, so $\tup{3, 0}$ and $\tup{1, 1}$ both belong to $\hat{\monrelM}(3)$.
    A least element of $\hat{\monrelM}(3)$ would have to be below both of them in the product order, hence of the form $\tup{\elnc{1}, 0}$ with $\elnc{1} \leq 1$.
    But such coins are worth at most one franc, so they are not in $\hat{\monrelM}(3)$.
    Thus $\hat{\monrelM}(3)$ has no least element, and \cref{lem:which-monrels-are-companions} shows that \monrelM is not a companion.

    Similarly, if $\posAset$ is non-empty, the empty relation $\posA \profto \posB$ is not a companion, since each $\hat{\monrelM}(\posAel)$ is empty.
    So monotone relations are \emph{strictly} more general than monotone functions, just as relations are strictly more general than functions.
\end{example}

\subsubsection{Companions and composition}

For sets, \cref{lem:comprelfun} says that composing graphs of functions as relations gives the graph of the composite function.
The analogous statement holds for companions and conjoints; for conjoints the order of composition is reversed, as for the transpose of relations.

\begin{lemma}
    \label{lem:companion-conjoint-composition}
    Let $\mapa \colon \posA \mto \posB$ and $\mapb \colon \posB \mto \posC$ be \SY{monotone functions} between \SY{preorders}.
    Then
    \begin{equation}
        \comp{\mapa \mthen \mapb} = \comp{\mapa} \mthen \comp{\mapb}
    \end{equation}
    and
    \begin{equation}
        \conj{\mapa \mthen \mapb} = \conj{\mapb} \mthen \conj{\mapa}.
    \end{equation}
\end{lemma}
\begin{proof}
    Recall that $(\mapa \mthen \mapb)(\posAel) = \mapb(\mapa(\posAel))$.

    Companions: by definition, $\inrel{\posAel}{(\comp{\mapa} \mthen \comp{\mapb})}{\posCel}$ holds if and only if there is $\posBel \setin \posBset$ with $\mapa(\posAel) \posBleq \posBel$ and $\mapb(\posBel) \posCleq \posCel$.
    If there is such a $\posBel$, then, since $\mapb$ is monotone, $\mapb(\mapa(\posAel)) \posCleq \mapb(\posBel) \posCleq \posCel$, so $\inrel{\posAel}{\comp{\mapa \mthen \mapb}}{\posCel}$ by transitivity.
    Conversely, if $\mapb(\mapa(\posAel)) \posCleq \posCel$, then $\posBel = \mapa(\posAel)$ works, since $\mapa(\posAel) \posBleq \mapa(\posAel)$ by reflexivity.

    Conjoints: by definition, $\inrel{\posCel}{(\conj{\mapb} \mthen \conj{\mapa})}{\posAel}$ holds if and only if there is $\posBel \setin \posBset$ with $\posCel \posCleq \mapb(\posBel)$ and $\posBel \posBleq \mapa(\posAel)$.
    If there is such a $\posBel$, then $\posCel \posCleq \mapb(\posBel) \posCleq \mapb(\mapa(\posAel))$, so $\inrel{\posCel}{\conj{\mapa \mthen \mapb}}{\posAel}$.
    Conversely, if $\posCel \posCleq \mapb(\mapa(\posAel))$, then $\posBel = \mapa(\posAel)$ works, again by reflexivity.
\end{proof}

Together with \cref{exa:companion-identity}, this says that taking companions is compatible with identities and with composition.
Later in the book, when we have the language of categories and functors, we will be able to express this more concisely.

\begin{example}[Course admission, revisited]
    In \cref{exa:monrel-courses}, the functions $\stylemaps{req} \colon \posA \mto \realswithleq$ and $\stylemaps{avg} \colon \posB \mto \realswithleq$ are monotone, by the way we defined the preorders on \posA and \posB.
    The monotone relation ``course $\ela$ is open to student $\elb$'' is the composite
    \begin{equation}
        \monrelM = \comp{\stylemaps{req}} \mthen \conj{\stylemaps{avg}} \colon \posA \profto \realswithleq \profto \posB.
    \end{equation}
    Indeed, $\inrel{\ela}{(\comp{\stylemaps{req}} \mthen \conj{\stylemaps{avg}})}{\elb}$ holds if and only if there is a number $\elc \setin \reals$ with $\stylemaps{req}(\ela) \leq \elc$ and $\elc \leq \stylemaps{avg}(\elb)$.
    If there is such a $\elc$, then $\stylemaps{req}(\ela) \leq \stylemaps{avg}(\elb)$ by transitivity; conversely, if $\stylemaps{req}(\ela) \leq \stylemaps{avg}(\elb)$, we can take $\elc = \stylemaps{req}(\ela)$.
    The intermediate preorder $\realswithleq$ is the ``grade scale'' through which courses and students are compared.
\end{example}

\subsection{Summary: monotone relations generalize monotone functions}

Let us collect the parallels between the passage from functions to relations (\cref{chap:relation}) and the passage from monotone functions to monotone relations.
\begin{itemize}
    \item Relations between sets are the special case of monotone relations between discrete preorders (\cref{exa:monrel-discrete}), just as functions between sets are the monotone functions between discrete preorders.
    \item A function is viewed as a relation through its graph (\cref{def:functions_as_relations}); a monotone function is viewed as a monotone relation through its companion (\cref{def:monrel-companion-conjoint}), which is the graph closed upward in the target (\cref{lem:companion-via-graph}).
          There is also a second way, the conjoint, which goes in the opposite direction.
    \item Just as a function can be recovered from its graph, a monotone function into a poset can be recovered from its companion (\cref{lem:companion-determines-function}).
          Not every monotone relation is a companion (\cref{lem:which-monrels-are-companions}), just as not every relation is a function.
    \item Relations compose, and composing graphs of functions gives the graph of the function composite (\cref{lem:comprelfun}); monotone relations compose (\cref{lem:monrel-composition}), and composing companions gives the companion of the composite (\cref{lem:companion-conjoint-composition}).
    \item The identity relation on a set is equality; the identity monotone relation on a preorder is the preorder relation itself (\cref{lem:monrel-leq-neutral}), which is the companion of the identity function (\cref{exa:companion-identity}).
    \item A relation can be encoded as a function $\setA \cartprod \setB \sto \boolset$; a monotone relation can be encoded as a monotone function $\posAop \Ptimes \posB \mto \Bool$ (\cref{lem:monrel-as-bool-map}).
    \item A relation can be described by the function $\hat\phi_\relA \colon \setA \sto \powerset(\setB)$ of \cref{rem:rel-three-fun-descriptions}; in the next chapter we will see that a design problem can similarly be described by a monotone function into a poset of upper sets, which answers all its queries at once (\cref{sec:dp-querying}).
\end{itemize}


